\documentclass[11pt,a4paper]{article}
\usepackage[utf8]{inputenc}
\usepackage{amsmath,amssymb,amsthm}
\usepackage{hyperref}
\usepackage{booktabs}
\usepackage{listings}
\usepackage[margin=2.5cm]{geometry}

\theoremstyle{definition}
\newtheorem{definition}{Definition}[section]
\newtheorem{theorem}{Theorem}[section]
\newtheorem{lemma}[theorem]{Lemma}
\newtheorem{corollary}[theorem]{Corollary}
\newtheorem{example}{Example}[section]

\lstdefinelanguage{MeTTa}{
  morekeywords={import,def,prob-rule,prob-goal,prob-assume,exec-conj,bdd-wmc,collapse,superpose,match,if,let,apply,bddVar},
  morestring=[b]",
  morecomment=[l]{;},
  sensitive=true
}
\lstset{
  basicstyle=\ttfamily\small,
  keywordstyle=\bfseries,
  commentstyle=\itshape,
  columns=flexible,
  breaklines=true
}

\title{ProbMeTTa: ProbLog in MeTTa\\[0.3em]
\large Formalization, Implementation, and Verified Correctness}

\author{}
\date{}

\begin{document}
\maketitle

\begin{abstract}
We present ProbMeTTa, an implementation of ProbLog's distribution semantics
in the MeTTa programming language, running on the PeTTa (Prolog-backed MeTTa)
runtime. We formalize ProbLog's distribution semantics in Lean~4 and prove
that ProbMeTTa's BDD-based inference correctly computes ProbLog query
probabilities. The formalization comprises approximately 7{,}200 lines of
Lean with zero unproven assumptions (\texttt{sorry}), covering BDD
construction correctness, Weighted Model Counting, compilation
soundness and completeness, an explicit runtime evaluation relation
for ProbMeTTa's source-level query operations, a normalized source surface
for the top-level ProbMeTTa operators, a literal annotated-disjunction
builder core for \texttt{new-id}, \texttt{build-ad-helper}, and
\texttt{build-ad-rules}, an LP-level literal matching layer for nonground
fallback, a literal state-space layer for \texttt{\&prob-facts}/
\texttt{\&prob-templates}, and an end-to-end theorem connecting the literal
runtime presentation back to the proved semantics. All 27 ProbMeTTa test cases
pass, and 7
representative programs are additionally verified by kernel-checked
conformance tests in Lean.
\end{abstract}

\section{Introduction}

ProbLog~\cite{deraedt2007} is a probabilistic extension of Prolog based on
Sato's distribution semantics~\cite{sato1995}. It combines the declarative
expressiveness of logic programming with principled probabilistic reasoning,
and has been applied to link prediction, bioinformatics, and probabilistic
databases.

MeTTa is a functional-relational programming language designed for
AI~reasoning. PeTTa is a Prolog-backed MeTTa implementation that
compiles MeTTa expressions to SWI-Prolog.

ProbMeTTa implements ProbLog's exact inference engine in MeTTa using Binary
Decision Diagrams (BDDs) for knowledge compilation and Weighted Model
Counting (WMC) for probability computation, following the approach of
Fierens et~al.~\cite{fierens2015}.

This paper presents:
\begin{enumerate}
\item A self-contained formalization of ProbLog's distribution semantics in Lean~4
\item ProbMeTTa: a working ProbLog implementation in MeTTa (27/27 tests passing)
\item A formal proof that ProbMeTTa's query/runtime core, source surface,
  literal matcher/state layers, and top-level presentation correctly compute
  ProbLog probabilities
\item Kernel-checked conformance tests linking the Lean formalization to the MeTTa code
\end{enumerate}

\section{ProbLog: Distribution Semantics}

\subsection{Programs}

A ProbLog program $P$ consists of:
\begin{itemize}
\item \textbf{Probabilistic facts}: $p_1\!::\!f_1,\ \ldots,\ p_n\!::\!f_n$ ---
  each ground atom $f_i$ is independently true with probability $p_i \in [0,1]$.
\item \textbf{Definite clauses}: $\mathit{head} \leftarrow \mathit{body}_1, \ldots, \mathit{body}_k$ ---
  standard Horn clauses.
\end{itemize}

\begin{example}[Alarm Network]
\begin{lstlisting}
0.1 :: burglary.
0.2 :: earthquake.
alarm :- burglary.
alarm :- earthquake.
query(alarm).
\end{lstlisting}
\end{example}

\subsection{Semantics}

ProbLog defines query probability via the \emph{distribution semantics}
\cite{sato1995}:

\begin{definition}[Total Choice]
A \emph{total choice} is an assignment $a : \{1,\ldots,n\} \to \{\texttt{true},\texttt{false}\}$
that independently sets each probabilistic fact to true or false.
\end{definition}

\begin{definition}[Residual Program]
For total choice $a$, the \emph{residual program} consists of the original
rules plus the probabilistic facts chosen true: $\mathit{KB}(a) = \mathit{rules} \cup \{f_i \mid a(i) = \texttt{true}\}$.
\end{definition}

\begin{definition}[Query Probability]
Query $q$ \emph{holds under $a$} iff $q$ is in the least Herbrand model of $\mathit{KB}(a)$.
The \emph{query probability} is:
\[
P(q) = \sum_{a} w(a) \cdot \mathbf{1}[q \text{ holds under } a]
\]
where $w(a) = \prod_{i=1}^{n} \bigl(a(i) \cdot p_i + (1-a(i)) \cdot (1-p_i)\bigr)$.
\end{definition}

For the alarm example: $P(\mathit{alarm}) = 1 - (1-0.1)(1-0.2) = 0.28$.

\section{BDD-Based Inference}

ProbLog computes query probabilities via knowledge compilation to Binary
Decision Diagrams (BDDs)~\cite{bryant1986}, following Fierens et~al.~\cite{fierens2015}.

\subsection{Binary Decision Diagrams}

A BDD over $n$ Boolean variables is a rooted directed acyclic graph with:
\begin{itemize}
\item Two terminal nodes: $\mathbf{0}$ (false) and $\mathbf{1}$ (true)
\item Internal nodes labeled by variable $v \in \{1,\ldots,n\}$, each with
  a \emph{lo} child (variable false) and \emph{hi} child (variable true)
\end{itemize}

A Reduced Ordered BDD (ROBDD) additionally satisfies:
\begin{enumerate}
\item \textbf{Ordered}: variables appear in consistent order on all paths
\item \textbf{Reduced}: no node has $\mathit{lo} = \mathit{hi}$ (elimination rule)
\end{enumerate}

\subsection{Weighted Model Counting}

Given a BDD $f$ and probability weights $p_1, \ldots, p_n$:
\begin{align}
\mathrm{WMC}(\mathbf{0}) &= 0 \\
\mathrm{WMC}(\mathbf{1}) &= 1 \\
\mathrm{WMC}(\mathrm{node}(v, \mathit{lo}, \mathit{hi})) &= (1-p_v) \cdot \mathrm{WMC}(\mathit{lo}) + p_v \cdot \mathrm{WMC}(\mathit{hi})
\end{align}

\begin{theorem}[WMC Correctness]
For any ROBDD $f$ representing Boolean function $\varphi$:
\[
\mathrm{WMC}(f) = \sum_{\substack{a \in \{0,1\}^n \\ \varphi(a) = 1}} \prod_{i=1}^{n} \bigl(a_i \cdot p_i + (1-a_i)(1-p_i)\bigr)
\]
\end{theorem}

This is exactly the distribution semantics probability when $\varphi$ is
the query's success predicate.

\subsection{Compilation}

ProbLog compiles queries to BDDs by proof enumeration:
\begin{enumerate}
\item Each probabilistic fact $f_i$ becomes BDD variable $i$
\item A conjunction $(b_1, \ldots, b_k)$ compiles to $\mathrm{AND}$ of sub-BDDs
\item Multiple derivations of the same head are $\mathrm{OR}$'d together
\end{enumerate}

\section{ProbMeTTa: Implementation in MeTTa}

ProbMeTTa implements the above algorithm in two MeTTa libraries:

\subsection{\texttt{lib\_bdd.metta}: BDD Operations}

\begin{lstlisting}[language=MeTTa]
;; BDD node: (bdd-node $id $var $lo $hi)
;; Terminals: bdd-0 (false), bdd-1 (true)

;; Make node with elimination rule
(= (mk $var $lo $hi)
    (if (== $lo $hi) $lo
        (case (once (match &bdd ...))
            ((Empty (let $id (new-id) ...))
             ($id $id)))))

;; Shannon expansion with memoization
(= (apply-bdd $op $f $g) ...)

;; Weighted Model Counting
(= (bdd-wmc $f $env)
    (if (is-terminal $f)
        (if (== $f bdd-1) 1.0 0.0)
        (let* (($var (bdd-var $f))
               ($val (env-lookup $var $env)))
            (+ (* (- 1 $val) (bdd-wmc (bdd-lo $f) $env))
               (* $val (bdd-wmc (bdd-hi $f) $env))))))
\end{lstlisting}

\subsection{\texttt{lib\_prob.metta}: Probabilistic Inference}

\begin{lstlisting}[language=MeTTa]
;; Declare probabilistic fact: p :: f
(= (:: $prob $goal)
    (progn (add-atom &prob-templates (prob-template $prob $goal))
           (add-atom &self (rule $goal (pf $goal)))))

;; Declare definite clause: head :- body
(= (prob-rule $body $head)
    (let $goals (body-goals $body)
        (add-atom &self (rule $head $goals))))

;; Resolve a goal, threading BDD trace
(= (prob-goal $goal $trace)
    (let $body (match &self (rule $goal $b) $b)
        (case $body
            (((pf $atom) (prob-assume $atom $trace))
             ($goals (exec-conj $goals $trace))))))

;; Top-level query
(= (?prob $goal)
    (let* (($final-bdd (?prob-bdd $goal))
           ($env (prob-wmc-env)))
        (bdd-wmc $final-bdd $env)))
\end{lstlisting}

\subsection{Test Results}

ProbMeTTa passes all 27 tests, covering:

\begin{table}[h]
\centering
\begin{tabular}{lll}
\toprule
\textbf{Test} & \textbf{ProbLog Feature} & \textbf{Expected} \\
\midrule
Alarm network & Noisy-OR, 3-hop chain & $P = 0.196$ \\
Fever chain & 6 rules, shared derivations & $P = 0.35168$ \\
Simple negation & Negation-as-failure & $P = 0.7$ \\
Conditioning & $P(A \mid B) = P(A \wedge B)/P(B)$ & $P = 0.357$ \\
Annotated disjunctions & Multi-valued random variables & $P = 0.6$ \\
Graph reachability & Recursive probabilistic rules & $P = 0.258$ \\
Overlapping rules & Multiple derivation paths & $P = 0.49$ \\
\bottomrule
\end{tabular}
\caption{Representative ProbMeTTa test results (all 27 pass).}
\end{table}

\section{Lean Formalization}

The formalization is in Lean~4 (v4.28.0 with Mathlib v4.28.0), totaling
approximately 7{,}200 lines in the verified ProbMeTTa slice with zero
\texttt{sorry} (unproven assumptions). The core compilation and WMC correctness
is approximately 1{,}700 lines; extensions cover annotated disjunctions,
first-order translation, the explicit runtime evaluation relation, the
normalized source surface, literal nonground matching, literal state-space and
presentation layers, and the end-to-end theorem surface.

\subsection{BDD Core (\texttt{Core.lean}, 169 lines)}

\begin{lstlisting}[language={}]
inductive BDD (n : N) where
  | zero : BDD n                                -- false
  | one  : BDD n                                -- true
  | node (v : Fin n) (lo hi : BDD n) : BDD n   -- if v then hi else lo

def BDD.eval (f : BDD n) (env : Fin n -> Bool) : Bool
\end{lstlisting}

The denotational semantics \texttt{BDD.eval} maps each BDD to the Boolean
function $(a_1,\ldots,a_n) \mapsto \{0,1\}$ it represents.

\subsection{WMC Correctness (\texttt{WMC.lean}, 239 lines)}

\begin{theorem}[\texttt{bdd\_wmc\_correct}]
For any well-formed BDD $f$ with probability weights $p_i \leq 1$:
\[
\texttt{bdd\_wmc}(f, \mathit{env}) = \sum_{a \in \{0,1\}^n} \mathbf{1}[f.\mathit{eval}(a)] \cdot \prod_{i=1}^n w(i, a_i)
\]
where $w(i, \texttt{true}) = p_i$ and $w(i, \texttt{false}) = 1 - p_i$.
\end{theorem}

The proof proceeds by structural induction on the BDD, with the key lemma
being the \emph{Shannon decomposition} for weighted sums: splitting the sum
by one variable's value and factoring out its weight.

\subsection{Operations Correctness (\texttt{Operations.lean}, 245 lines)}

\begin{theorem}[\texttt{apply\_eval}]
\[
(\texttt{apply}\ \mathit{op}\ f\ g).\mathit{eval}(a) = \mathit{op}(f.\mathit{eval}(a),\ g.\mathit{eval}(a))
\]
\end{theorem}

This proves that BDD construction via Shannon expansion preserves Boolean
semantics. Seven conformance tests verify the Lean BDD operations against
ProbMeTTa's expected results:

\begin{table}[h]
\centering
\begin{tabular}{lll}
\toprule
\textbf{Program} & \textbf{Boolean function} & \textbf{Verification} \\
\midrule
Alarm & $v_0 \vee v_1$ & \texttt{by simp} \\
Conjunction & $v_0 \wedge v_1$ & \texttt{by simp} \\
Negation & $\neg v_0$ & \texttt{by simp} \\
Neg.~conjunction & $v_0 \wedge \neg v_1$ & \texttt{by simp} \\
Overlapping & $(v_0 \wedge v_1) \vee v_2$ & \texttt{by simp} \\
Fever chain & $(v_2 \wedge v_3) \vee ((v_0 \vee v_1) \wedge v_2 \wedge v_4)$ & \texttt{by simp} \\
Calls-mary & $(v_0 \vee v_1) \wedge v_2$ & \texttt{by simp} \\
\bottomrule
\end{tabular}
\caption{Kernel-checked conformance tests. Each \texttt{by simp} is verified by Lean's kernel.}
\end{table}

\subsection{Compilation Correctness (\texttt{Compilation.lean}, 268 lines)}

The compilation relation \texttt{GroundBDDCompile} mirrors ProbMeTTa's
\texttt{prob-goal} procedure:

\begin{center}
\begin{tabular}{lll}
\toprule
\textbf{Constructor} & \textbf{ProbMeTTa function} & \textbf{Meaning} \\
\midrule
\texttt{.fact} & \texttt{prob-assume} & Probabilistic fact $\to$ BDD variable \\
\texttt{.ruleG} & \texttt{prob-goal + exec-conj} & Rule application $\to$ AND of body BDDs \\
\texttt{.disj} & \texttt{bdd-disjunction} & Multiple derivations $\to$ OR \\
\bottomrule
\end{tabular}
\end{center}

\begin{theorem}[Soundness: \texttt{GroundBDDCompile\_sound}]
If \texttt{GroundBDDCompile prog q f} and $f.\mathit{eval}(a) = \texttt{true}$,
then $q \in \mathrm{LHM}(\mathit{KB}(a))$.
\end{theorem}

\begin{theorem}[Completeness: \texttt{GroundBDDCompile\_complete}]
If $q \in \mathrm{LHM}(\mathit{KB}(a))$, then there exists $f$ such that
\texttt{GroundBDDCompile prog q f} and $f.\mathit{eval}(a) = \texttt{true}$.
\end{theorem}

Completeness is proved by induction on the $T_P$ iterate level, following
the characterization $\mathrm{LHM}(\mathit{KB}) = \bigcup_k T_P^k(\emptyset)$.
At each level, EDB facts compile to variable nodes, and rule applications
compile to AND/OR combinations of body BDDs obtained from the inductive
hypothesis.

\subsection{Runtime Core (\texttt{ProbMeTTaRuntimeCore.lean}, 393 lines)}

The formalization includes an explicit big-step evaluation relation
\texttt{ProbMeTTaCoreEval} for the actual source-level function names
from \texttt{lib\_prob.metta}:

\begin{center}
\begin{tabular}{ll}
\toprule
\textbf{Source function} & \textbf{Formalized operation} \\
\midrule
\texttt{prob-assume} & BDD variable from probabilistic fact \\
\texttt{prob-not} & Complete BDD + negation \\
\texttt{prob-goal} & Dispatch: positive $\to$ ground compile, negative $\to$ \texttt{prob-not} \\
\texttt{exec-conj} & AND-fold over goal list \\
\texttt{?prob-bdd} & OR over all derivations \\
\texttt{?prob-bdd-ev} & AND-fold over evidence-query BDDs \\
\bottomrule
\end{tabular}
\end{center}

\begin{theorem}[\texttt{probMeTTaCoreEval\_queryBDD\_iff\_ground}]
The runtime evaluation relation for \texttt{?prob-bdd} is equivalent to
the abstract compilation relation: both produce exactly the same BDD
for a given query.
\end{theorem}

\begin{theorem}[\texttt{exists\_ordered\_probMeTTa\_queryBDD\_exact}]
For any query $q$, there exists an ordered BDD $f$ such that
\texttt{ProbMeTTaCoreEval prog (.queryBDD q) f} and $f$ is
sound and complete for $q$.
\end{theorem}

This layer bridges the gap between the abstract compilation relation
(\texttt{GroundBDDCompile}) and the actual source-level runtime operations
in ProbMeTTa's \texttt{lib\_prob.metta}.

\subsection{Normalized Source Surface (\texttt{ProbMeTTaSourceSurface.lean}, 627 lines)}

On top of the runtime/query core, the formalization now includes a normalized
source surface for the actual top-level ProbMeTTa operators:
\texttt{prob-rule}, \texttt{=>}, \texttt{::}, \texttt{::adr}, \texttt{::=>},
\texttt{prob-wmc-env}, \texttt{?prob-bdd-ev}, \texttt{?prob}, and
\texttt{?prob-given}.

The normalization is explicit and intentional:
\begin{itemize}
\item rule bodies are stored as explicit lists of goal literals;
\item top-level probabilities are exact \texttt{ENNReal} values rather than
  the source runtime's decimal \texttt{round-to 5} presentation;
\item literal nonground fallback is handled separately by an LP-level matching
  layer, rather than being hidden inside a witness abstraction;
\item \texttt{?prob-bdd-ev} and \texttt{?prob-given} are formalized as explicit
  source operators rather than only as downstream semantic consequences.
\end{itemize}

\subsection{Literal Matching (\texttt{ProbMeTTaLiteralMatching.lean}, 239 lines)}

The source nonground branch
\verb|match &self (rule $goal $_) $goal|
is now formalized against the existing LP matching layer. The formal surface
returns the list of matching ground heads and the corresponding list of exact
\texttt{?prob}/\texttt{?prob-given} values, which is closer to the source
\texttt{match} nondeterminism than the earlier matched-head witness packaging.

\subsection{AD Runtime Core (\texttt{ProbMeTTaADRuntimeCore.lean}, 415 lines)}

The source-level annotated-disjunction builders in \texttt{lib\_prob.metta}
are also formalized explicitly:
\texttt{new-id}, \texttt{build-ad-helper}, \texttt{build-ad-rules},
\texttt{::adr}, and \texttt{::=>}.

This layer keeps the boundary honest by separating:
\begin{itemize}
\item a \emph{literal} builder layer, which threads fresh ids, cumulative mass,
  and the source runtime's last-head optimization;
\item the existing \emph{normalized} weighted-AD layer, which connects directly
  to the already-proved AD/WMC and conditioning theorems.
\end{itemize}

\subsection{Crown Theorem Bridge (\texttt{ProbMeTTaBridge.lean}, 349 lines)}

The crown theorem file composes all BDD results into the top-level equivalence:

\begin{theorem}[\texttt{problog\_full\_ground\_equivalence}]
For any ground stratified normal ProbLog program, there exist ordered BDDs
for query and evidence goals such that their WMC ratio equals the
distribution semantics conditional probability $P(Q \mid E)$.
\end{theorem}

The composition chain:
\begin{enumerate}
\item \textbf{Runtime exactness}: \texttt{probmetta\_runtime\_query\_bdd\_exact}
      proves the source-level \texttt{?prob-bdd} fragment has an ordered BDD
      that is true exactly on worlds where the query holds.
\item \textbf{WMC ratio identity}: \texttt{wmc\_ratio\_identity} proves that
      the ratio of two BDD WMC values equals the ratio of their weighted
      satisfiability sums.
\item \textbf{Conditioning bridge}: \texttt{goal\_conditioning\_semantic\_equivalence}
      connects $P(Q \wedge E) / P(E)$ to compiled goal BDDs.
\item \textbf{Evidence positivity}: \texttt{normal\_conditional\_probability}
      ensures the division is meaningful when $P(E) \neq 0$.
\end{enumerate}

This layer was implicit in the proof architecture diagram but is now
explicitly documented as a verified bridge (0~\texttt{sorry}).

\subsection{Literal State and Presentation Layers}

The remaining source-level runtime structure is now explicit too:
\begin{itemize}
\item \texttt{ProbMeTTaStateSpace.lean} (270 lines) formalizes a literal state
  for \texttt{\&prob-facts}, \texttt{\&prob-templates},
  \texttt{clear-prob!}, and \texttt{prob-wmc-env};
\item \texttt{ProbMeTTaPresentation.lean} (134 lines) formalizes the top-level
  \texttt{round-to 5} presentation as an explicit source output layer;
\item \texttt{ProbMeTTaEndToEnd.lean} (200 lines) states the crown end-to-end
  theorems linking the literal state/runtime/presentation layers back to the
  proved exact BDD/source semantics.
\end{itemize}

\noindent
This still does \emph{not} claim byte-for-byte replay of MeTTa reduction order.
What is proved is semantic agreement between the literal layers we modeled and
the already-proved exact ProbLog/BDD semantics.

\subsection{Proof Architecture}

The complete verified chain:

\[
\underbrace{\text{ProbLog program}}_{\texttt{ProbLogProgram}}
\xrightarrow{\text{literal state/source}}
\underbrace{\text{state + source surface}}_{\texttt{ProbMeTTaStateSpace},\,\texttt{ProbMeTTaSourceProgram}}
\xrightarrow{\text{runtime}}
\underbrace{\text{runtime query core}}_{\texttt{ProbMeTTaCoreEval}}
\overset{\cong}{\longleftrightarrow}
\underbrace{\text{BDD}}_{\texttt{GroundBDDCompile}}
\xrightarrow{\text{WMC}}
\underbrace{P(q)}_{\texttt{bdd\_wmc\_correct}}
\xrightarrow{\text{presentation}}
\underbrace{\texttt{round-to 5}}_{\texttt{ProbMeTTaPresentation}}
\]

The runtime relation \texttt{ProbMeTTaCoreEval} for
\texttt{?prob-bdd}/\texttt{?prob-bdd-ev} is proved equivalent to the abstract
BDD compilation layer, the normalized source surface packages the actual
top-level ProbMeTTa operators on top of that bridge, the literal matching/state
layers reconnect the source runtime branches that depend on \texttt{match} and
mutable spaces, and the presentation layer closes the loop at the displayed
top-level output.

\begin{itemize}
\item \textbf{Soundness + Completeness}: the compiled BDD evaluates to true
  under assignment $a$ if and only if the query holds in the corresponding
  ProbLog world.
\item \textbf{WMC Correctness}: Weighted Model Counting on the BDD equals
  the distribution semantics weighted sum.
\item \textbf{Therefore}: $\texttt{bdd\_wmc}(f, \mathit{env}) = P(q)$.
\end{itemize}

\section{WM-PLN Subsumes ProbLog}

The World Model -- Probabilistic Logic Networks (WM-PLN) framework
subsumes ProbLog: it computes the same probabilities and additionally
provides confidence values and incremental evidence revision.

\subsection{Subsumption: PLN Fuzzy-OR = ProbLog Noisy-OR}

PLN's fuzzy-OR formula for independent causes is:
\[
s_1 + s_2 - s_1 \cdot s_2 = 1 - (1-s_1)(1-s_2)
\]

This is algebraically identical to the noisy-OR formula that arises from
ProbLog's distribution semantics for OR-pattern programs (proved in Lean:
\texttt{fuzzyOrMulti\_eq\_noisyOrMulti}, 0~\texttt{sorry}).

For the alarm example, both ProbMeTTa and PLN give $P(\mathit{alarm}) = 0.28$:

\begin{lstlisting}[language=MeTTa]
;; ProbMeTTa (BDD-based ProbLog):
!(?prob (alarm))                   ;; => 0.28

;; PLN fuzzy-OR (native, no BDDs):
!(pln-fuzzy-or 0.1 0.2)           ;; => 0.28
\end{lstlisting}

\subsection{Extension: Evidence-Based Reasoning with Confidence}

ProbLog stores fixed-point probabilities. PLN stores \emph{evidence}:
pairs of positive and negative observation counts $(n^+, n^-)$.

\begin{definition}[Evidence and Simple Truth Value]
An evidence pair $(n^+, n^-)$ yields:
\begin{align}
\mathrm{strength} &= \frac{n^+}{n^+ + n^-} \\
\mathrm{confidence} &= \frac{n^+ + n^-}{n^+ + n^- + \kappa}
\end{align}
where $\kappa$ is a prior context parameter.
\end{definition}

Strength matches ProbLog's probability. Confidence measures how much
evidence supports the estimate --- something ProbLog cannot express.

\subsection{Evidence Revision: What ProbLog Cannot Do}

ProbLog stores $P(\mathit{burglary}) = 0.4$ as a weight.
It cannot distinguish whether this came from 4 observations out of 10
(weak evidence) or 400 out of 1{,}000 (strong evidence).
When new data arrives, ProbLog requires external relearning.

PLN stores the counts $(n^+, n^-)$ and revises incrementally via
\emph{evidence aggregation} (\texttt{evidence-hplus}):
\[
(n^+_1, n^-_1) \oplus (n^+_2, n^-_2) = (n^+_1 + n^+_2,\ n^-_1 + n^-_2)
\]

\subsection{Demonstration: Convergence to Known Truth}

We construct a synthetic benchmark with known ground truth
($P(\mathit{alarm}) = 0.28$) and reveal observations in three batches:

\begin{table}[h]
\centering
\begin{tabular}{lcccc}
\toprule
\textbf{After} & \textbf{P(alarm)} & \textbf{Error} & \textbf{Confidence} & \textbf{Method} \\
\midrule
10 windows  & 0.4000 & 0.1200 & 0.50 & ProbLog (stuck) / WM-PLN \\
40 windows  & 0.3025 & 0.0225 & 0.80 & WM-PLN revision \\
100 windows & 0.2791 & 0.0009 & 0.91 & WM-PLN revision \\
\midrule
100 (rebuilt) & 0.2791 & 0.0009 & --- & ProbLog (external relearning) \\
\bottomrule
\end{tabular}
\caption{Evidence revision converges toward truth with 133$\times$ error reduction.
ProbLog matches the final answer only after a complete program rebuild.}
\end{table}

The estimate converges from 0.40 (initial small-sample error) to 0.2791
(within 0.001 of truth), while confidence rises from 0.50 to 0.91.
ProbLog can match the final answer, but only via an external
weight-relearning step --- it has no internal revision mechanism.

The key insight: ProbLog represents $P = 0.4$ and cannot distinguish
$4/10$ from $400/1000$. Those two cases should revise very differently
when new data arrives, but ProbLog represents them identically.
WM-PLN keeps the sufficient statistics $(n^+, n^-)$, so
the amount of revision is principled.

\section{Getting ProbMeTTa Running on PeTTa}

ProbMeTTa was originally developed for Hyperon Experimental MeTTa. Running
it on PeTTa required fixing four compatibility issues:

\begin{enumerate}
\item \textbf{\texttt{is-ground} undefined}: PeTTa lacks this built-in.
  Fixed by wrapping SWI-Prolog's \texttt{ground/1} via a one-line Prolog
  predicate.

\item \textbf{Singleton list patterns}: PeTTa's \texttt{(\$h . \$t)} cons
  pattern does not match singletons \texttt{(x)}. Fixed by using
  \texttt{car-atom}/\texttt{cdr-atom} instead of dotted-pair patterns.

\item \textbf{\texttt{case} branch ordering}: PeTTa's \texttt{case} with
  \texttt{once(match ...)} requires the \texttt{Empty} branch first.

\item \textbf{\texttt{=>} not evaluated}: PeTTa treats \texttt{=>} as
  structural, not as a reducible function. Fixed by introducing
  \texttt{prob-rule} as the rule declaration entrypoint.
\end{enumerate}

After these fixes, all 27 ProbMeTTa tests pass on PeTTa.

\section{Conclusion}

We have shown three things:

\begin{enumerate}
\item \textbf{ProbMeTTa correctly implements ProbLog.}
A formal proof in Lean~4 (approximately 7{,}200 lines, 0~\texttt{sorry}) covers BDD
construction, Weighted Model Counting, compilation soundness and
completeness, an explicit runtime evaluation relation for
ProbMeTTa's source-level query operations, a normalized source surface for the
top-level operators, a literal nonground matching layer, literal state and
presentation layers, and a literal annotated-disjunction builder core.
Kernel-checked conformance tests link the abstract proof to the concrete MeTTa
implementation. All 27 ProbMeTTa tests pass.

\item \textbf{WM-PLN subsumes ProbLog.}
PLN's fuzzy-OR is algebraically identical to ProbLog's noisy-OR
(proved: \texttt{fuzzyOrMulti\_eq\_noisyOrMulti}). For any ground definite
clause ProbLog program, WM-PLN computes the same query probability
(proved via BDD compilation soundness and completeness; negation-as-failure
and annotated disjunctions additionally verified empirically via 27/27
ProbMeTTa tests).

\item \textbf{WM-PLN extends ProbLog with confidence and revision.}
By storing evidence counts rather than point probabilities, WM-PLN
provides confidence values and supports incremental revision.
In a synthetic benchmark, evidence revision reduces estimation error
by 133$\times$ (from 0.12 to 0.0009) while confidence rises from 0.50
to 0.91. ProbLog can match the final answer only via external relearning.
\end{enumerate}

The key conceptual point: ProbLog discards the sufficient statistics
behind its probability estimates. A weight of $0.4$ could represent
$4/10$ or $400/1000$ observations --- ProbLog cannot tell the difference.
WM-PLN retains the evidence $(n^+, n^-)$, making revision principled
and incremental.

\subsection*{Availability}

\begin{itemize}
\item ProbMeTTa: \url{https://github.com/zariuq/ProbMeTTa}
\item PeTTa: \url{https://github.com/zariuq/PeTTa} (branch \texttt{wmPLN})
\item Lean formalization: \texttt{Mettapedia/Logic/BDD/} in the Mettapedia project
\item ProbLog (upstream): \url{https://github.com/ML-KULeuven/problog}
\end{itemize}

\begin{thebibliography}{9}

\bibitem{deraedt2007}
L.~De~Raedt, A.~Kimmig, H.~Toivonen.
\newblock ProbLog: A Probabilistic Prolog and its Application in Link Discovery.
\newblock \emph{IJCAI}, 2007.

\bibitem{sato1995}
T.~Sato.
\newblock A Statistical Learning Method for Logic Programs with Distribution Semantics.
\newblock \emph{ICLP}, 1995.

\bibitem{fierens2015}
D.~Fierens, G.~Van~den~Broeck, J.~Renkens, D.~Shterionov, B.~Gutmann,
I.~Thon, G.~Janssens, L.~De~Raedt.
\newblock Inference and Learning in Probabilistic Logic Programs using
Weighted Boolean Formulas.
\newblock \emph{Theory and Practice of Logic Programming}, 15(3), 2015.

\bibitem{bryant1986}
R.~E.~Bryant.
\newblock Graph-Based Algorithms for Boolean Function Manipulation.
\newblock \emph{IEEE Transactions on Computers}, C-35(8), 1986.

\bibitem{vanemden1976}
M.~H.~van~Emden, R.~A.~Kowalski.
\newblock The Semantics of Predicate Logic as a Programming Language.
\newblock \emph{Journal of the ACM}, 23(4), 1976.

\end{thebibliography}

\end{document}
